\scanpage{278}
% archive foliation: 277

l'Ecole de Mathématique \unsure{est de} 6 heures hebdomadaires, est dans ses attributions d'être le conseiller pédagogique des jeunes assistants et maître-assistants.

5°/ Il appartient à chacun des Maîtres de Conférences et Professeurs mathématiciens de la Faculté des Sciences d'Orsay en fonctions avant la rentrée d'Octobre 1968 de décider de son affectation au Laboratoire de Recherche ou à l'Ecole de Mathématiques, à titre permanent. \uline{Il est entendu qu'en cas de déséquilibre provisoire de l'ensemble des deux sections au profit du Laboratoire de Recherche, la vocation spéciale des membres de celui-ci ne pourra être respectée que dans la mesure où elle n'entre pas en conflit avec les impératifs de l'enseignement à assurer.}

$\overline{\underline{\mathrm{II}}}$ Critères de nomination.

La qualification nécessaire pour être nommé Maître de Conférence dans l'une ou l'autre section du Département de Mathématiques est le titre de Docteur ès Sciences Mathématiques, mention ``Enseignement'' pour les candidats à l'Ecole de Mathématique, mention ``Recherche'' pour les candidats au Laboratoire de Recherche Mathématique.

$\overline{\underline{\mathrm{III}}}$ Recrutement.

1°/ La nomination à tout poste d'enseignant ou de chercheur dans l'une ou l'autre section du département de
